%% PACKAGES
\usepackage[utf8]{inputenc}
\usepackage[T1]{fontenc}
\usepackage[french]{babel}

\usepackage{minted}
\setminted[ocaml]{
    %bgcolor = black!3!white,
    frame=leftline, framesep=6pt, rulecolor=expli,
    linenos=true, numbersep=4pt, xleftmargin=20pt,
    breaklines=true,
    tabsize=2,
}

\usepackage{amsmath, amsfonts, amssymb}
\usepackage{textcomp, lmodern}
\usepackage{cancel}

\usepackage{array, multirow}
\usepackage{tabularx}

\usepackage{tikz}
\usetikzlibrary{shapes}  % pour écrire \node[ellipse] par ex
\usetikzlibrary{positioning}  % pour écrire \node[below rigth = 3pt and 5pt]
\usetikzlibrary{decorations.pathreplacing}  % pour les accolades
\usetikzlibrary{patterns}  % pour les hachures

\usepackage{caption}  % les personaliser (les centrer par ex)
\usepackage{changepage}  % pour élargir la page avec adjustwidth
\usepackage{calc}  % pour pouvoir écrire \textwidth-1cm
\usepackage{xspace}  % espaces après les commandes texte


\input{couleurs.tex}

%% PARAMÈTRES BEAMER
\usetheme{Montpellier}
\usecolortheme{seahorse}

% marges
\setbeamersize{text margin left=0.7cm}
\setbeamersize{text margin right=0.7cm}

% en-têtes et bas de page
\setbeamertemplate{navigation symbols}{}
\setbeamertemplate{footline}[frame number]
\setbeamertemplate{headline}{
  % la premiere ligne
  \begin{beamercolorbox}[ht=0.42cm, vmode]{section in head/foot}
    \hspace{0.4cm} \insertshorttitle
    \hspace*{0.1cm}- \initiales - {\insertshortdate}
    \vspace*{0.08cm}
  \end{beamercolorbox}
  % la deuxième ligne
  \begin{beamercolorbox}[ht=0.4cm, vmode]{subsection in head/foot}
    % titre de la section si elle n’est pas 0
    \ifnum\value{section}=0{}
    \else{ \hspace{0.8cm} \thesection - \insertsectionhead }
    \fi
    % séparateur + titre de la sous-section si elle n’est pas 0
    \ifnum\value{subsection}=0{}
    \else{
        \hspace*{0.1cm} \couleur{$\bullet$} \hspace*{0.1cm}
        \thesection.\thesubsection \, \insertsubsectionhead
    }
    \fi
    \vspace*{0.12cm}
  \end{beamercolorbox}
  % \vspace*{-0.03cm} % pour pas qu'il n’y ait d'espace avec la ligne de
                      % frametitle
}

\setbeamertemplate{frametitle}{
    \vspace*{-0.04cm}
    \begin{beamercolorbox}[ht=0.8cm, wd=\paperwidth, vmode]{frametitle}
        \hspace{0.3cm} \insertframetitle \vspace*{0.1cm}
    \end{beamercolorbox}
}

% commande pour ajuster l'alignement vertical des titres sans lettres
% descendantes
\newcommand{\esp}{\\[-0.5cm]} % version qui marche avec le package minted


% couleurs
\setbeamercolor{structure}{fg=turquoiseFonce!70!black}

\setbeamercolor{block title}{fg=turquoiseFonce!70!black,bg=vertdEau}
\setbeamercolor{block body}{bg=vertdEau!10!white}

\setbeamercolor{block title alerted}{bg=vertdEau!85!white,fg=turquoiseFonce!80!black}
\setbeamercolor{block body alerted}{bg=vertdEau!8!white}

\setbeamercolor{block title example}{bg=vertdEau,fg=turquoiseFonce}
\setbeamercolor{block body example}{bg=vertdEau!10!white}
\setbeamercolor{example text}{fg=blue!20!turquoise}

% TOC
\setbeamertemplate{section in toc}[sections numbered]

% -----------------------------------------------|
% Plan qui s'affiche au début de chaque section %|
\AtBeginSection[]{                              %|
    \begin{frame}[plain]                        %|
        \frametitle{Plan\\[0.1cm]}              %|
        \tableofcontents[                       %|
            currentsection,                     %|
            hideothersubsections,               %|
            subsubsectionstyle=hide]            %|
        \addtocounter{framenumber}{-1}          %|
    \end{frame}                                 %|
}                                               %|
% -----------------------------------------------|


% commande pour les ref sur les slides
\newcommand{\bandeauREF}[1]{
    \noindent\makebox[\textwidth][l]{%
        \hspace{-\dimexpr\oddsidemargin+1in}%
            \colorbox{expli!20!white}{%
                \parbox{\dimexpr\paperwidth-2\fboxsep}{
                    \footnotesize\textcolor{expli!80!black}{#1}
                }
            }
    }
}
